\documentclass[11pt]{article}
\usepackage[a4paper,margin=1in]{geometry}
\usepackage{amsmath,amssymb,amsthm}
\usepackage{graphicx}
\usepackage[hidelinks]{hyperref}

\newtheorem{theorem}{Theorem}
\newtheorem{lemma}{Lemma}
\newtheorem{corollary}{Corollary}
\newtheorem{remark}{Remark}

\DeclareMathOperator{\Lip}{Lip}

\title{Separable Domains Give All-Codomain Projections in \(\Lip_0(X,Y)\)}
\author{Automatic math-discovery packet}
\date{June 29, 2026}

\begin{document}
\maketitle

\section*{Source question}

The source paper is Anil Kumar Karn and Arindam Mandal,
\emph{Position of \(L(X,Y)\) in \(\Lip_0(X,Y)\)}, arXiv:2506.09324.
Question 1.1 asks whether \(L(X,Y)\) is complemented in \(\Lip_0(X,Y)\).

\begin{figure}[h]
\centering
\includegraphics[width=0.93\textwidth]{figures/open_problem_crop.png}
\caption{Question 1.1 asks whether \(L(X,Y)\) is complemented in \(\Lip_0(X,Y)\).}
\end{figure}

\clearpage

\section*{Source inputs}

Karn--Mandal define \(F_Y(X)\) as the closed linear span of the evaluation
maps \(\delta_x^Y:\Lip_0(X,Y)\to Y\), and define a contractive linear
barycenter map
\[
  \beta_X^Y:F_Y(X)\longrightarrow X,\qquad \beta_X^Y(\delta_x^Y)=x.
\]
They prove that
\[
  \Delta_X^Y:L(F_Y(X),Y)\longrightarrow \Lip_0(X,Y),
  \qquad \Delta_X^Y(A)=A\circ\delta_X^Y,
\]
is a surjective linear isometry, and that
\[
  L(X,Y)={}^{\diamondsuit}\ker(\beta_X^Y).
\]
Equivalently, under \(\Delta_X^Y\), the subspace \(L(X,Y)\) corresponds to
the operators \(A:F_Y(X)\to Y\) that vanish on \(\ker(\beta_X^Y)\), i.e. to
the operators factoring through \(\beta_X^Y\).

We also use the Godefroy--Kalton lifting theorem: if \(X\) is separable, then
the usual scalar Lipschitz-free barycenter map
\[
  \beta_X:\mathcal F(X)\longrightarrow X
\]
has a norm-one bounded linear right inverse. We cite this in the form
recorded in Ferenczi--Kaufmann--Pernecka, \emph{Equivariant liftings in
Lipschitz-free spaces}, arXiv:2501.06984, Theorem 2.

\section*{Proof intuition}

The source paper's generalized space \(F_Y(X)\) looks \(Y\)-valued, but its
underlying molecule norm is actually the scalar Lipschitz-free norm. Testing
a \(Y\)-valued Lipschitz function against \(Y^*\) gives the upper bound by
the scalar free norm, while scalar Lipschitz functions multiplied by a fixed
unit vector of \(Y\) give the reverse inequality.

Thus any linear lifting of the scalar barycenter map
\(\beta_X:\mathcal F(X)\to X\) becomes a linear lifting of
\(\beta_X^Y:F_Y(X)\to X\). A linear lifting splits
\[
  0\longrightarrow \ker(\beta_X^Y)\longrightarrow F_Y(X)
  \stackrel{\beta_X^Y}{\longrightarrow} X\longrightarrow 0,
\]
and this split gives a projection on operator spaces. Transferring it through
\(\Delta_X^Y\) gives a projection from \(\Lip_0(X,Y)\) onto \(L(X,Y)\).

\section*{The scalar free space is already present}

\begin{lemma}
Let \(Y\neq\{0\}\). The linear map initially defined on molecules by
\[
  U\left(\sum_{i=1}^n a_i\delta_{x_i}\right)
  =
  \sum_{i=1}^n a_i\delta_{x_i}^Y
\]
extends to a surjective linear isometry
\[
  U:\mathcal F(X)\longrightarrow F_Y(X).
\]
Moreover \(\beta_X^Y U=\beta_X\).
\end{lemma}

\begin{proof}
For a finite molecule \(m=\sum_i a_i\delta_{x_i}\), its \(F_Y(X)\)-norm is
\[
  \|Um\|_{F_Y(X)}
  =
  \sup_{\Lip(f)\leq 1}
  \left\|\sum_i a_i f(x_i)\right\|,
\]
where \(f\) ranges over \(\Lip_0(X,Y)\). For every \(y^*\in B_{Y^*}\),
\(y^*\circ f\in B_{\Lip_0(X)}\). Hence
\[
  \left\|\sum_i a_i f(x_i)\right\|
  =
  \sup_{y^*\in B_{Y^*}}
  \left|\sum_i a_i\,y^*f(x_i)\right|
  \leq \|m\|_{\mathcal F(X)}.
\]
Thus \(\|Um\|_{F_Y(X)}\leq \|m\|_{\mathcal F(X)}\).

Conversely, choose \(y_0\in Y\) with \(\|y_0\|=1\). For every
\(g\in B_{\Lip_0(X)}\), the map \(f_g(x)=g(x)y_0\) belongs to
\(B_{\Lip_0(X,Y)}\), and
\[
  \left\|\sum_i a_i f_g(x_i)\right\|
  =
  \left|\sum_i a_i g(x_i)\right|.
\]
Taking the supremum over \(g\) gives
\(\|Um\|_{F_Y(X)}\geq \|m\|_{\mathcal F(X)}\). Therefore \(U\) is an
isometry on the molecule span, and it extends to an isometry
\(\mathcal F(X)\to F_Y(X)\). Its range is dense by definition of \(F_Y(X)\),
and closed because it is isometric, so it is onto. Finally,
\(\beta_X^Y U(\delta_x)=\beta_X^Y(\delta_x^Y)=x=\beta_X(\delta_x)\), and
the identity extends by continuity.
\end{proof}

\section*{Lifting-property projection}

\begin{theorem}
Suppose the scalar barycenter map \(\beta_X:\mathcal F(X)\to X\) has a
bounded linear right inverse \(S:X\to\mathcal F(X)\), with
\(\beta_XS=I_X\). Then for every Banach space \(Y\), the subspace \(L(X,Y)\)
is complemented in \(\Lip_0(X,Y)\). More precisely, if \(Y\neq\{0\}\), there
is a projection of norm at most \(\|S\|\); if \(Y=\{0\}\), the assertion is
trivial.
\end{theorem}

\begin{proof}
Assume \(Y\neq\{0\}\). Let \(U:\mathcal F(X)\to F_Y(X)\) be the isometry of
the lemma and put \(S_Y=US:X\to F_Y(X)\). Then
\[
  \beta_X^Y S_Y=\beta_X^YUS=\beta_XS=I_X.
\]
Define
\[
  Q:L(F_Y(X),Y)\longrightarrow L(F_Y(X),Y),
  \qquad Q(A)=A S_Y\beta_X^Y.
\]
Since \(\|\beta_X^Y\|\leq 1\), we have \(\|Q\|\leq\|S_Y\|\leq\|S\|\).
Also \(Q^2=Q\), because \(\beta_X^YS_Y=I_X\).

The range of \(Q\) consists of operators that factor through \(\beta_X^Y\),
hence vanish on \(\ker(\beta_X^Y)\). Conversely, if
\(A\in L(F_Y(X),Y)\) vanishes on \(\ker(\beta_X^Y)\), then \(A\) factors
through \(\beta_X^Y\), say \(A=B\beta_X^Y\) for some \(B\in L(X,Y)\), and
\[
  Q(A)=B\beta_X^Y S_Y\beta_X^Y=B\beta_X^Y=A.
\]
Thus \(Q\) is a projection onto the operators corresponding, under the
Karn--Mandal isometry \(\Delta_X^Y\), to \(L(X,Y)\). Therefore
\[
  \Pi=\Delta_X^Y Q(\Delta_X^Y)^{-1}
\]
is a bounded projection from \(\Lip_0(X,Y)\) onto \(L(X,Y)\), with
\(\|\Pi\|\leq\|S\|\).
\end{proof}

\begin{corollary}
If \(X\) is separable, then for every Banach space \(Y\), \(L(X,Y)\) is
one-complemented in \(\Lip_0(X,Y)\).
\end{corollary}

\begin{proof}
By the Godefroy--Kalton lifting theorem, separable \(X\) has a norm-one
linear right inverse \(S:X\to\mathcal F(X)\) for \(\beta_X\). Apply the
theorem.
\end{proof}

\section*{Limitations and novelty check}

This is not a full answer to Question 1.1 for arbitrary \(X\) and \(Y\). The
proof uses the scalar lifting property of \(X\), and Ferenczi--Kaufmann--
Pernecka recall that some Banach spaces fail this lifting property, citing
Godefroy--Kalton's examples. Thus this route cannot simply be pushed to all
domains.

The result is nevertheless independent of the previous codomain-duality
route: it assumes nothing about \(Y\). Cheap run-index searches found only
the previous arXiv:2506.09324 bidual-complemented-codomain packet and an
adjacent two-Lipschitz tensor-splitting packet, but no exact packet for the
one-variable \(\Lip_0(X,Y)\) separable-domain/all-codomain theorem.

\section*{Human-review recommendation}

Review as a likely valid partial result. The key points to check are:
first, the norm equality between \(F_Y(X)\) and the scalar
\(\mathcal F(X)\); second, that the projection \(A\mapsto AS_Y\beta_X^Y\)
has range exactly the operators vanishing on \(\ker(\beta_X^Y)\), which the
source paper identifies with \(L(X,Y)\) after applying \(\Delta_X^Y\).

\begin{thebibliography}{9}
\bibitem{KarnMandal}
A. K. Karn and A. Mandal,
\emph{Position of \(L(X,Y)\) in \(\Lip_0(X,Y)\)},
arXiv:2506.09324, 2025.

\bibitem{FerencziKaufmannPernecka}
V. Ferenczi, P. Kaufmann, and E. Pernecka,
\emph{Equivariant liftings in Lipschitz-free spaces},
arXiv:2501.06984, 2025.

\bibitem{GodefroyKalton}
G. Godefroy and N. J. Kalton,
\emph{Lipschitz-free Banach spaces},
Studia Math. 159 (2003), 121--141.
\end{thebibliography}

\end{document}
